\section{A quantitative reduction of the ramified global cutoff problem}
\label{sec:ACcutoffbridge}

This section continues the attempt to prove T-CUT from AC-T11.
The calculations below use the actual theta images and their
orthogonal projections.  They distinguish the scalar prolate
problem from the additional arithmetic norm.
Fix a primitive nonprincipal Dirichlet character \(\chi\) of
conductor \(q\), with \(\chi(-1)=(-1)^\epsilon\).
For \(\Lambda>1\), let \(\psi_1,\ldots,\psi_N\) be the prolate
vectors of parity \(\epsilon\) chosen in AC-T11, supported in
\([-\Lambda,\Lambda]\), with the conductor-adjusted frequency
parameter \(\Lambda/q\).  The intended full family has indices
up to \(4\Lambda^2/q\).
Put
\[
 e_j(x)=E_\chi\psi_j(x)
   =\sqrt{x}\sum_{n\geq1}\chi(n)\psi_j(nx),\quad
 G_{ij}=\langle e_i,e_j\rangle,\quad
 S_\Lambda=1_{[\Lambda^{-1},\Lambda]}.
\]
Inner products are conjugate-linear in the first variable.
The Hilbert space is \(L^2(\mathbb R_{>0},dx/x)\).

\begin{proposition}[AC-T21: the exact arithmetic Gram weight]
\label{AC-T21}
Define
\[
 m_j(\tau)=\int_0^\infty
       \psi_j(x)x^{-1/2+i\tau}\,dx,\qquad
 m(\tau)=(m_1(\tau),\ldots,m_N(\tau)).
\]
Then
\[
 G=\frac1{2\pi}\int_{\mathbb R}
      |L(1/2+i\tau,\chi)|^2 m(\tau)^*m(\tau)\,d\tau,
 \tag{TB.1}
\]
and, for the orthogonal projection \(Q_\Lambda\) onto
\(\operatorname{span}\{e_j\}\),
\[
 \operatorname{Tr}(Q_\Lambda V(h))
   =\frac1{2\pi}\int_{\mathbb R}\widehat h_+(\tau)
       K_{\chi,\Lambda}(\tau)\,d\tau,\qquad
 K_{\chi,\Lambda}(\tau)=
 |L(1/2+i\tau,\chi)|^2
           m(\tau)G^{-1}m(\tau)^* .
 \tag{TB.2}
\]
These integrals are absolutely convergent for bounded
\(\widehat h_+\).

Conjugating the original time-frequency rectangle by dilation
\(\sqrt q\) replaces the scalar prolate parameter by
\(\Lambda/\sqrt q\).  This dilation multiplies every \(m_j(\tau)\)
by the same phase \(q^{i\tau/2}\), and so leaves (TB.1)--(TB.2)
unchanged apart from the replacement of the scalar prolate basis.
The weight \(|L(1/2+i\tau,\chi)|^2\) remains.
For nonprincipal \(\chi\), it is not a constant multiple of
\(|\zeta(1/2+i\tau)|^2\).
\end{proposition}
\begin{proof}
Let \(B_\chi=\sup_M|\sum_{n\leq M}\chi(n)|<\infty\).
Summation by parts bounds
\(\sum\chi(n)\psi_j(nx)\) uniformly in \(x>0\) by
\(B_\chi\) times the variation of the zero extension of
\(\psi_j\) on the positive half-line.
It also vanishes for \(x>\Lambda\).
Thus \(e_j\) lies in \(L^1(dx/x)\cap L^2(dx/x)\);
near zero it is \(O(\sqrt x)\).

For \(\operatorname{Re}s>1\), absolute convergence gives
\[
 \int_0^\infty\left(\sum_{n\geq1}\chi(n)\psi_j(nx)\right)
                      x^{s-1}\,dx
       =L(s,\chi)\int_0^\infty\psi_j(x)x^{s-1}\,dx.
\]
The summation-by-parts bound continues the left side
holomorphically to \(\operatorname{Re}s>0\); the right side
has the same continuation.  At \(s=1/2+i\tau\) this is the
Fourier transform of \(e_j(e^t)\).
Plancherel proves (TB.1).  The finite-rank projection formula
\(Q_\Lambda=EG^{-1}E^*\) from AC-T11 gives (TB.2).
The integral of \(K_{\chi,\Lambda}\) divided by \(2\pi\) is
\(\operatorname{tr}(G^{-1}G)=N\), which proves the claimed
absolute convergence.

If \(\psi(x)=q^{-1/4}\phi(x/\sqrt q)\), its Mellin
transform at \(1/2+i\tau\) is \(q^{i\tau/2}\) times that of
\(\phi\).  The common phases cancel in both matrix formulas.
Finally suppose the two stated weights were proportional
with constant \(c>0\).  Analytic continuation from the line
\(\operatorname{Re}s=1/2\) would imply
\[
 L(s,\chi)L(1-s,\overline\chi)
       =c\,\zeta(s)\zeta(1-s).
\]
The left side is entire.  The right side has a pole at \(s=1\),
because \(\zeta(0)=-1/2\).  This is impossible.
\end{proof}

\begin{proposition}[AC-T22: exact tail and principal-angle errors]
\label{AC-T22}
Let
\[
 H=E^*(1-S_\Lambda)E,\qquad
 A_\Lambda=G^{-1/2}HG^{-1/2}.
\]
Then \(0\leq A_\Lambda<1\).
If \(Q_\Lambda^0\) denotes the orthogonal projection onto
\(\operatorname{span}\{S_\Lambda e_j\}\), then
\[
 \|Q_\Lambda-Q_\Lambda^0\|_2^2
       =2\operatorname{tr}A_\Lambda,\qquad
 \|Q_\Lambda-Q_\Lambda^0\|_1
       =2\operatorname{tr}\sqrt{A_\Lambda}
       \leq2\sqrt{N\operatorname{tr}A_\Lambda}.
 \tag{TB.3}
\]
In particular
\[
 \left|\operatorname{Tr}\bigl((Q_\Lambda-Q_\Lambda^0)V(h)\bigr)\right|
       \leq 2\|h\|_{L^1(dx/x)}
                     \operatorname{tr}\sqrt{A_\Lambda}.
 \tag{TB.4}
\]
The tail matrix has the exact conductor-sensitive expression
\[
 H_{ij}=\int_{\Lambda/q}^\infty
   \overline{E_{\overline\chi}(\widehat\psi_i)(y)}
             E_{\overline\chi}(\widehat\psi_j)(y)\,\frac{dy}{y}.
 \tag{TB.5}
\]
For piecewise smooth cutoff vectors the theta sums on the
right can equivalently be understood through the Poisson
identity almost everywhere.
\end{proposition}
\begin{proof}
Every \(e_j\) vanishes above \(\Lambda\), so \(H\) is
precisely the Gram matrix of the tails below \(\Lambda^{-1}\).
Finite Fourier transformation of \(\chi\), followed by ordinary
Poisson summation, gives
\[
 E_\chi\psi(x)=\omega_\chi
       E_{\overline\chi}\widehat\psi(1/(qx)),
 \qquad |\omega_\chi|=1.
 \tag{TB.6}
\]
The phase is the normalized Gauss sum, with the harmless
sign determined by the Fourier convention.  There are no
constant terms because \(\chi\) is nonprincipal.
For compactly supported piecewise smooth functions this is
the ordinary Poisson identity for a periodized function of
bounded variation, away from its countable endpoint set.
Changing \(y=1/(qx)\) in the tail integral proves (TB.5).

The vectors \(S_\Lambda e_j\) are linearly independent.
Indeed on
\((\max(\Lambda^{-1},\Lambda/2),\Lambda)\) the theta sum
has only its \(n=1\) term.  If a linear combination of the
truncated images vanishes there, the same combination of
prolate functions vanishes on that open interval.
These functions are analytic inside \((-\Lambda,\Lambda)\),
so the combination is identically zero.  Their original
linear independence gives the assertion.
Thus \(G-H\) is positive definite and \(A_\Lambda<1\).

Put \(Z=EG^{-1/2}\), an isometry from \(\mathbb C^N\).
Then \(Z^*S_\Lambda Z=1-A_\Lambda\) and
\[
 Q_\Lambda^0=S_\Lambda Z(1-A_\Lambda)^{-1}Z^*S_\Lambda.
\]
Diagonalize \(A_\Lambda\) with eigenvalues \(\alpha_j\).
The pairs consisting of \(Z u_j\) and
\(S_\Lambda Zu_j/\sqrt{1-\alpha_j}\) are orthogonal for
different \(j\), with mutual inner product
\(\sqrt{1-\alpha_j}\) in the \(j\)-th pair.
On their two-dimensional span the difference of the
rank-one projections has eigenvalues
\(\pm\sqrt{\alpha_j}\).  Zero angles contribute nothing.
Summation proves both norm identities in (TB.3), and
Cauchy--Schwarz proves its final inequality.
Finally \(\|V(h)\|\leq\|h\|_1\) gives (TB.4).
\end{proof}

\begin{proposition}[AC-T23: a computable sufficient trace gate]
\label{AC-T23}
Write \(L=\log\Lambda\) and let
\(k^0_{\chi,\Lambda}(\tau)\) be the Fourier diagonal density
of \(Q_\Lambda^0\), defined from any orthonormal basis of its
range.  Then
\[
 0\leq k^0_{\chi,\Lambda}(\tau)\leq2L,\qquad
 \operatorname{Tr}\bigl((S_\Lambda-Q_\Lambda^0)V(h)\bigr)
 =\frac1{2\pi}\int_{\mathbb R}
       \widehat h_+(\tau)\bigl(2L-k^0_{\chi,\Lambda}(\tau)\bigr)
       \,d\tau.
 \tag{TB.7}
\]
Consequently the following two estimates suffice for the
nonprincipal-character assertion T-CUT:
\[
 \operatorname{tr}\sqrt{A_\Lambda}=o(1),\qquad
 \frac{(2L-k^0_{\chi,\Lambda}(\tau))\,d\tau}{2\pi}
  \ \longrightarrow\ \sum_\gamma m_{\chi,\gamma}\delta_\gamma
 \tag{TB.8}
\]
where the second convergence is tested against
\(\widehat h_+\) for every \(h\in C_c^\infty(\mathbb R_{>0})\).
The second estimate includes control at high frequency,
not just vague convergence on compact sets.

A direct variation estimate gives the fully explicit bound
\[
 \operatorname{tr}A_\Lambda
 \leq\frac{B_\chi^2}{\Lambda\,\lambda_{\min}(G)}
        \sum_{j=1}^N
          \operatorname{Var}_{(0,\infty)}(\psi_j)^2 .
 \tag{TB.9}
\]
This estimate alone does not establish (TB.8) for the
growing prolate family of size \(N\asymp\Lambda^2/q\).
\end{proposition}
\begin{proof}
In logarithmic coordinates let \(z_1,\ldots,z_N\) be an
orthonormal basis of the range of \(Q_\Lambda^0\),
supported in \([-L,L]\).
Then
\[
 k^0_{\chi,\Lambda}(\tau)
       =\sum_{j=1}^N\left|
                  \int_{-L}^Lz_j(t)e^{i\tau t}\,dt\right|^2.
 \]
Bessel's inequality against the vector
\(t\mapsto e^{-i\tau t}\) in \(L^2[-L,L]\) gives the upper
bound \(2L\).  Fourier inversion and the finite-rank trace
give (TB.7); the interval convolution trace is
\(2Lh(1)\), as in AC-T19.
Combining (TB.7) and (TB.4) proves sufficiency of (TB.8).

For (TB.9), summation by parts as in AC-T21 gives
\[
 |e_j(x)|\leq B_\chi\sqrt{x}\,
                     \operatorname{Var}_{(0,\infty)}(\psi_j).
 \]
Integration on \(0<x<\Lambda^{-1}\) bounds \(H_{jj}\)
by \(B_\chi^2\operatorname{Var}(\psi_j)^2/\Lambda\).
Since \(G^{-1}\leq\lambda_{\min}(G)^{-1}\), tracing gives
(TB.9).  No positive lower bound for
\(\lambda_{\min}(G)\), uniform through the prolate threshold,
has been proved here; nor does (TB.9) control the second
limit in (TB.8).
\end{proof}

The second limit in (TB.8) is the continuous counterpart of the
finite-zero computation in C99, Lemma VIII.3, equations
(28)--(29).  In the function-field proof only finitely many
zeros occur in the fundamental frequency domain, and a dimension
count identifies the image exactly.
For a number-field character there are infinitely many zeros.
Under \(A\), compactly supported logarithmic functions cannot vanish at
all of them unless zero, by AC-T13's Jensen argument.
The finite image is therefore not characterized by imposing
all those exact vanishing conditions.
Equations (TB.1), (TB.5), and (TB.8) identify the additional
arithmetic estimates needed after the scalar prolate rescaling.
They do not assert that these estimates are false, or that they
are logically independent of \(A\).

\subsection{An exact closure consequence and why its cutoff is different}

\begin{corollary}[AC-T24: one-sided closures under \(A\)]
\label{AC-T24}
In the preceding nonprincipal character sector, \(A\) gives
\[
 \overline{E_\chi(C_c^\infty(-R,R)^\epsilon)}
       =L^2((0,R),dx/x),\qquad R>0.
 \tag{TB.10}
\]
Here the superscript fixes parity.  It follows that
\[
 \overline{E_\chi(C_c^\infty(-\Lambda,\Lambda)^\epsilon)}
 \ \cap\
 \overline{E_\chi\bigl(\mathcal F^{-1}
              C_c^\infty(-\Lambda/q,\Lambda/q)^\epsilon\bigr)}
       =L^2([\Lambda^{-1},\Lambda],dx/x).
 \tag{TB.11}
\]
\end{corollary}
\begin{proof}
Burnol, \emph{An adelic causality problem related to abelian
\(L\)-functions}, arXiv:\path{math/0001013v3},
Theorem 3.2 and its number-field proof on pp.~13--15,
identifies this one-sided closure with the Hardy subspace
multiplied by the Blaschke product of the zeros to the right
of the critical line.  Burnol's multiplier \(V\) is identically
one on every nonprincipal character sector (display on
printed p.~13), so no extra operator modifies this closure.
The fixed finite character vector and
arbitrary smooth archimedean factors used there suffice.
Under \(A\) that product is one, proving (TB.10) at \(R=1\);
dilation gives arbitrary \(R\).
Apply (TB.10) to \(\overline\chi\) and use (TB.6) to obtain
the second closed space in (TB.11) as
\(L^2((\Lambda^{-1},\infty),dx/x)\).
Intersecting the two support spaces proves (TB.11).
\end{proof}

The completed intersection in (TB.11) is the whole norm window.
Its projection is \(S_\Lambda\), whose subtracted trace defect
is identically zero.  The prolate-image projection is the
finite-rank \(Q_\Lambda\); its proposed defect contains the
critical-zero measure.  Thus replacing the finite approximation
by the intersection of the completed one-sided ranges would
discard the spectral information that T-CUT seeks to recover.
The exact closure result is useful and is a consequence of \(A\);
the rate and trace comparison in (TB.8) remain separate work.
